\section{Related Work}
\label{sec:related}

\subsection{Static Factor Weighting}

Traditional multi-factor models use fixed weights. Fama and French~\cite{fama1993} equal-weighted their factor portfolios. Barra~\cite{rosenberg1974} and MSCI derive weights through risk model optimization.

Gu et al.~\cite{gu2020} found neural networks outperform linear methods for return prediction, but focused on accuracy rather than adaptive weighting. Green et al.~\cite{green2017} documented IC variation over time, motivating dynamic weighting without implementing one.

\subsection{Factor Timing}

Arnott et al.~\cite{arnott2016timing} showed factor premiums exhibit predictable variation. Cohen et al.~\cite{cohen2020bounds} established theoretical bounds: perfect timing is impossible, but meaningful timing is achievable.

Blitz et al.~\cite{blitz2020factormomentum} documented factor momentum: factors that performed well recently tend to continue. This directly supports IC-based weighting. Kelly and Jiang~\cite{kelly2023characteristics} interpreted firm characteristics as covariance risk proxies.

\subsection{Adaptive Weighting}

IC-based weighting uses rolling Spearman correlation to measure factor effectiveness. Prior work~\cite{qian2015,grinold2000} used IC to identify crowding and adjust exposures, but without systematic walk-forward validation.

Risk model optimization~\cite{clarke2016} adjusts weights based on volatility forecasts but does not incorporate IC dynamics. Neural network approaches~\cite{gu2020} achieve high in-sample fit but lack interpretability.

\subsection{Regime-Dependent Effects}

Asness et al.~\cite{asness2013} showed momentum returns are positive in trending markets but negative during reversals. Novy-Marx~\cite{novy-marx2013} documented quality factor effectiveness during downturns. In Chinese markets, Liu et al.~\cite{liu2019} identified distinct factor behaviors. Liu and Ma~\cite{liu2022review} noted higher volatility and more frequent regime transitions.

\subsection{Walk-Forward Validation}

Walk-forward optimization addresses overfitting by separating training and testing periods. Beggs and Chi~\cite{beggs2018} showed it provides more realistic estimates than traditional backtesting. This methodology has been applied to strategy development~\cite{arnold2019} and risk validation~\cite{prado2018}.

\subsection{Positioning}

We differ from existing work in five ways. First, we implement real-time IC tracking with exponential decay rather than fixed weights. Second, we preserve interpretability rather than using black-box neural networks. Third, we provide walk-forward validation on Chinese A-shares across multiple regime transitions including the 2022 correction. Fourth, we frame IC-based weighting as lightweight factor timing with theoretical support. Fifth, we integrate risk controls that address practical deployment concerns~\cite{demiguel2009,jorion1996,rockafellar2000}.
